%*********************************************************************
%	DISCUSSION
%*********************************************************************
\chapter{Discussion}
\label{ch:discussion}
\comment{What Did We Learn}
%---------------------------------------------------------------------
There are several results that are discussed in this report,this chapter aims to provide a summary of these different results with further regard to the overarching project and surrounding literature
\section{Optimal Configuration}
It was seen that primary factor regarding the cameras capture configuration was resolution. It can be seen in \todo{Reference the results} that a higher resolution produces much better accuracy in the 3D reconstruction. The default operation of the camera uses a resolution of 848$\times$480, which performs around \todo{Insert metric difference here} worse than a higher resolution of 1280$\times$720. And it can also be seen that while filter and presets do have some effect on the accuracy, there is a minimal change \todo{use actual metrics here}, especially since the capture program records the raw depth stream, any filter alterations can be applied post capture. This can help to optimise the capture time, by aleviating some of the processing requirements of the host machine.

It can also be seen that capture distance has a large effect on the performance of the system, and that generally the closer the camera is to the surface, the better accuracy can be obtained. This however does inversely correspond with the desired operation of the overarching project, since a lower scan distance, would reduce the available scan region for analysis. furthermore, the D456 camera is optimised for long range outdoor scans, and as such has a higher minimum operating distance of only 600mm. This limits how close the camera can be placed to the surface, since operating at this boundary can lead some tilit and rotations dipping below this boundary. It was generally seen that during initial experiments, that at operatingheights of 687mm began to introduce more significant depth anomolies, due toi the below minimum depth, at high tilits such as 17\textdegree. It can be seen that there is a noticible change in accuracy \todo{Reference the results of 3D characterisaiton} as the operating height increases. \todo{Reference a paper that finds similar things, from bibliography}. For this reason it is generally recommended to operate at a height of around 700\,mm for highest accuracy, and then use the maximum height for the largest possible scan region. And it is generally recommended to have a reference scan at both heights, so that a comparison can be made.

It can be seen from angle analysis, that higher tilt angles  \todo{reference the appendix plots for accuracy vs tilt of 3D objects}, that accuracy actually increases at higher tilit angles, this is likey due to the camera able to capture more information about the highly sloped features of the reference objects, wjhen tilted. This contradicts the results of planar characterisation, which generally show an inverse trend with plane accuracy across height. However where the two results agree, is that the completnetess of the 3D object, does generally degrade at the heigher angles. This is attributed ot the fact that at higher angles, the peaks of the reference objects tend to hide the valeys at these high angles.

\section{Fusion and Roughness Estimation}

The results of fusion show that fusion provides good accuracy reconstruction, when compared to single scans. Fusion is able ro significantly reduce the noise present in in scans \todo{Reference the visual images}. Showing that it is closer in terms of surface representation than a single scan. One important distincion here is that results from rugosity evaluation, show that a single scan results in a rugosity metric that is closed to the reference objects true rugosity. Whilst at face value this seems as if a single scan is better able to represent the surface rugosity, there is a major detail that changes the findings of this result. The reconstruction capability of the D456, is unable to completely match the profile of the reference objects, especially in the small egg box surface case, both for single and fused scans. Why this matters is that fundamentally, any scan of the reference surface will have a lower rugosity than the true reference object. How this changes the results, is the fact that the single scans have more noise present in them than the fused scans. This arteficially increases the roughness of the surface, thus arteficially increasing the rugosity, to closer match the reference object. Further evidence of this can be seen when inspecting the scans of the sand surface. When presented with a flat surface, the single scan produces a metric with a higher rugosity than that of the fused scan. The flat surface should have a rugosity that is close to 1, being a predominantly flat surface. This shows that the fused scan is able to reduce the small depth noise present in single scans, and thus more accurately matches the surface state at any point.

Overall this characterisation finds that the biggest improvements that are made with respect to best accuracy is to minimise the height of the depth camera, and to utilise multi-view scan fusion. \todo{Give a nice summary that ties in different research, and ties into the overarching project}



%*********************************************************************